\documentclass[11pt]{article}
% Shared preamble for the printable mini lessons (1.4 onward).
%   \documentclass[11pt]{article}
%   % Shared preamble for the printable mini lessons (1.4 onward).
%   \documentclass[11pt]{article}
%   % Shared preamble for the printable mini lessons (1.4 onward).
%   \documentclass[11pt]{article}
%   \input{../preamble}
% Build the key with:  pdflatex -jobname=<name>_KEY "\def\iskey{}\input{<name>}"
% Figures shared with the web page live in figures/Unit_X/<name>.tex and are
% pulled in with \fig{Unit_X/<name>}; the standalone package drops their own
% preamble, so their macros come from figures/preamble.tex, loaded here.
\input{../../figures/preamble}
\usepackage[margin=0.75in]{geometry}
\usepackage{array}
\usepackage{enumitem}
\usepackage{amsmath}
\usepackage[helvet]{sfmath}  % math in the same sans font as the text
\usepackage{standalone}
\pagestyle{empty}
\setlength{\parindent}{0pt}
\setlist[enumerate,1]{leftmargin=1.5em,itemsep=1em}
\setlist[enumerate,2]{label=(\alph*),leftmargin=2em,itemsep=2.2em}

\newif\ifkey
\ifdefined\iskey \keytrue \fi
\colorlet{keyred}{red!75!black}
\newcommand{\ans}[1]{\ifkey\leavevmode{\color{keyred}#1}\fi}
% vertical work space on the student copy, less on the key
\newcommand{\work}[1]{\ifkey\vspace{0.6em}\else\vspace{#1}\fi}
\newcommand{\fig}[1]{\input{../../figures/#1}}
% Title line: \lessontitle{Quartiles, IQR, and Box Plots}
\newcommand{\lessontitle}[1]{{\Large\bfseries #1\ifkey\ \textcolor{keyred}{KEY}\fi}\hfill Name \underline{\hspace{2in}}\par\vspace{0.6em}}
\newcommand{\partner}{\par\vspace{0.4em}Work with a partner. Compare answers when you finish each problem.\par}

% Build the key with:  pdflatex -jobname=<name>_KEY "\def\iskey{}\input{<name>}"
% Figures shared with the web page live in figures/Unit_X/<name>.tex and are
% pulled in with \fig{Unit_X/<name>}; the standalone package drops their own
% preamble, so their macros come from figures/preamble.tex, loaded here.
% Shared preamble for web figures.  Each figure is a standalone .tex that
% \input's this file; scripts/make-figures.sh turns it into an SVG.
\usepackage[T1]{fontenc}
\usepackage{helvet}
\renewcommand{\familydefault}{\sfdefault}
\usepackage{tikz}
\usetikzlibrary{decorations.pathreplacing}
% axis numbers in the text font, not math italic
\pgfkeys{/pgf/number format/.cd,assume math mode=true,1000 sep={}}

% dot plot: \dotplot{min}{max}{value/count,...}{axis label}
\newcommand{\dotplot}[4]{%
\begin{tikzpicture}[x=0.9cm,y=0.33cm]
  \draw[thick] (#1-0.5,0) -- (#2+0.5,0);
  \foreach \t in {#1,...,#2} {\draw (\t,0.25) -- (\t,-0.25) node[below] {\t};}
  \foreach \v/\c in {#3} {\ifnum\c>0 \foreach \k in {1,...,\c} {\fill (\v,\k*1.15) circle (0.16cm);}\fi}
  \node at ({(#1+#2)/2},-2.3) {#4};
\end{tikzpicture}}

% shape sketch histogram (12 bins): \shapehist{heights}{label}
\newcommand{\shapehist}[2]{%
  \foreach \h [count=\i from 0] in {#1} \draw[fill=gray!30] (\i,0) rectangle (\i+1,\h);
  \draw (-0.3,0) -- (12.3,0);
  \node[font=\bfseries] at (6,-2.6) {#2};}
\newcommand{\meanmark}[1]{\fill (#1,0) -- ++(-0.3,-0.75) -- ++(0.6,0) -- cycle;}
\newcommand{\medmark}[2]{\draw[very thick,dashed] (#1,0) -- (#1,#2);}

% number line with minor ticks: \nline{min}{max}{label step}{axis label}
\newcommand{\nline}[4]{%
  \draw[thick] (#1-0.5,0) -- (#2+0.5,0);
  \foreach \t in {#1,...,#2} \draw (\t,0.12) -- (\t,-0.12);
  \pgfmathtruncatemacro{\nxt}{#1+#3}
  \foreach \t in {#1,\nxt,...,#2} \draw (\t,0.2) -- (\t,-0.2) node[below] {\t};
  \node at ({(#1+#2)/2},-0.95) {#4};}
% box plot above the line: \bplot{min}{Q1}{med}{Q3}{max}
\newcommand{\bplot}[5]{%
  \draw[thick] (#1,1.0) -- (#2,1.0);
  \draw[thick] (#4,1.0) -- (#5,1.0);
  \draw[thick] (#1,0.75) -- (#1,1.25);
  \draw[thick] (#5,0.75) -- (#5,1.25);
  \draw[thick,fill=gray!20] (#2,0.55) rectangle (#4,1.45);
  \draw[very thick] (#3,0.55) -- (#3,1.45);}

% ---- added for 1.4 onward ----
% axis on y=0: \axis{min}{max}{label step}{tick step}{axis label}
\newcommand{\axis}[5]{%
  \draw[thick] (#1,0) -- (#2,0);
  \pgfmathsetmacro{\axT}{#1+#4}\foreach \t in {#1,\axT,...,#2} \draw (\t,0.1) -- (\t,-0.1);
  \pgfmathsetmacro{\axL}{#1+#3}\foreach \t in {#1,\axL,...,#2}
    \draw (\t,0.18) -- (\t,-0.18) node[below] {\pgfmathprintnumber{\t}};
  \node at ({(#1+#2)/2},-0.95) {#5};}
% box plot centred at height y: \bplotat{y}{min}{Q1}{med}{Q3}{max}
\newcommand{\bplotat}[6]{%
  \draw[thick] (#2,#1) -- (#3,#1);
  \draw[thick] (#5,#1) -- (#6,#1);
  \draw[thick] (#2,#1-0.22) -- (#2,#1+0.22);
  \draw[thick] (#6,#1-0.22) -- (#6,#1+0.22);
  \draw[thick,fill=gray!20] (#3,#1-0.42) rectangle (#5,#1+0.42);
  \draw[very thick] (#4,#1-0.42) -- (#4,#1+0.42);}
% an outlier dot: \outlier{x}{y}
\newcommand{\outlier}[2]{\fill (#1,#2) circle (2.4pt);}

\usepackage[margin=0.75in]{geometry}
\usepackage{array}
\usepackage{enumitem}
\usepackage{amsmath}
\usepackage[helvet]{sfmath}  % math in the same sans font as the text
\usepackage{standalone}
\pagestyle{empty}
\setlength{\parindent}{0pt}
\setlist[enumerate,1]{leftmargin=1.5em,itemsep=1em}
\setlist[enumerate,2]{label=(\alph*),leftmargin=2em,itemsep=2.2em}

\newif\ifkey
\ifdefined\iskey \keytrue \fi
\colorlet{keyred}{red!75!black}
\newcommand{\ans}[1]{\ifkey\leavevmode{\color{keyred}#1}\fi}
% vertical work space on the student copy, less on the key
\newcommand{\work}[1]{\ifkey\vspace{0.6em}\else\vspace{#1}\fi}
\newcommand{\fig}[1]{\input{../../figures/#1}}
% Title line: \lessontitle{Quartiles, IQR, and Box Plots}
\newcommand{\lessontitle}[1]{{\Large\bfseries #1\ifkey\ \textcolor{keyred}{KEY}\fi}\hfill Name \underline{\hspace{2in}}\par\vspace{0.6em}}
\newcommand{\partner}{\par\vspace{0.4em}Work with a partner. Compare answers when you finish each problem.\par}

% Build the key with:  pdflatex -jobname=<name>_KEY "\def\iskey{}\input{<name>}"
% Figures shared with the web page live in figures/Unit_X/<name>.tex and are
% pulled in with \fig{Unit_X/<name>}; the standalone package drops their own
% preamble, so their macros come from figures/preamble.tex, loaded here.
% Shared preamble for web figures.  Each figure is a standalone .tex that
% \input's this file; scripts/make-figures.sh turns it into an SVG.
\usepackage[T1]{fontenc}
\usepackage{helvet}
\renewcommand{\familydefault}{\sfdefault}
\usepackage{tikz}
\usetikzlibrary{decorations.pathreplacing}
% axis numbers in the text font, not math italic
\pgfkeys{/pgf/number format/.cd,assume math mode=true,1000 sep={}}

% dot plot: \dotplot{min}{max}{value/count,...}{axis label}
\newcommand{\dotplot}[4]{%
\begin{tikzpicture}[x=0.9cm,y=0.33cm]
  \draw[thick] (#1-0.5,0) -- (#2+0.5,0);
  \foreach \t in {#1,...,#2} {\draw (\t,0.25) -- (\t,-0.25) node[below] {\t};}
  \foreach \v/\c in {#3} {\ifnum\c>0 \foreach \k in {1,...,\c} {\fill (\v,\k*1.15) circle (0.16cm);}\fi}
  \node at ({(#1+#2)/2},-2.3) {#4};
\end{tikzpicture}}

% shape sketch histogram (12 bins): \shapehist{heights}{label}
\newcommand{\shapehist}[2]{%
  \foreach \h [count=\i from 0] in {#1} \draw[fill=gray!30] (\i,0) rectangle (\i+1,\h);
  \draw (-0.3,0) -- (12.3,0);
  \node[font=\bfseries] at (6,-2.6) {#2};}
\newcommand{\meanmark}[1]{\fill (#1,0) -- ++(-0.3,-0.75) -- ++(0.6,0) -- cycle;}
\newcommand{\medmark}[2]{\draw[very thick,dashed] (#1,0) -- (#1,#2);}

% number line with minor ticks: \nline{min}{max}{label step}{axis label}
\newcommand{\nline}[4]{%
  \draw[thick] (#1-0.5,0) -- (#2+0.5,0);
  \foreach \t in {#1,...,#2} \draw (\t,0.12) -- (\t,-0.12);
  \pgfmathtruncatemacro{\nxt}{#1+#3}
  \foreach \t in {#1,\nxt,...,#2} \draw (\t,0.2) -- (\t,-0.2) node[below] {\t};
  \node at ({(#1+#2)/2},-0.95) {#4};}
% box plot above the line: \bplot{min}{Q1}{med}{Q3}{max}
\newcommand{\bplot}[5]{%
  \draw[thick] (#1,1.0) -- (#2,1.0);
  \draw[thick] (#4,1.0) -- (#5,1.0);
  \draw[thick] (#1,0.75) -- (#1,1.25);
  \draw[thick] (#5,0.75) -- (#5,1.25);
  \draw[thick,fill=gray!20] (#2,0.55) rectangle (#4,1.45);
  \draw[very thick] (#3,0.55) -- (#3,1.45);}

% ---- added for 1.4 onward ----
% axis on y=0: \axis{min}{max}{label step}{tick step}{axis label}
\newcommand{\axis}[5]{%
  \draw[thick] (#1,0) -- (#2,0);
  \pgfmathsetmacro{\axT}{#1+#4}\foreach \t in {#1,\axT,...,#2} \draw (\t,0.1) -- (\t,-0.1);
  \pgfmathsetmacro{\axL}{#1+#3}\foreach \t in {#1,\axL,...,#2}
    \draw (\t,0.18) -- (\t,-0.18) node[below] {\pgfmathprintnumber{\t}};
  \node at ({(#1+#2)/2},-0.95) {#5};}
% box plot centred at height y: \bplotat{y}{min}{Q1}{med}{Q3}{max}
\newcommand{\bplotat}[6]{%
  \draw[thick] (#2,#1) -- (#3,#1);
  \draw[thick] (#5,#1) -- (#6,#1);
  \draw[thick] (#2,#1-0.22) -- (#2,#1+0.22);
  \draw[thick] (#6,#1-0.22) -- (#6,#1+0.22);
  \draw[thick,fill=gray!20] (#3,#1-0.42) rectangle (#5,#1+0.42);
  \draw[very thick] (#4,#1-0.42) -- (#4,#1+0.42);}
% an outlier dot: \outlier{x}{y}
\newcommand{\outlier}[2]{\fill (#1,#2) circle (2.4pt);}

\usepackage[margin=0.75in]{geometry}
\usepackage{array}
\usepackage{enumitem}
\usepackage{amsmath}
\usepackage[helvet]{sfmath}  % math in the same sans font as the text
\usepackage{standalone}
\pagestyle{empty}
\setlength{\parindent}{0pt}
\setlist[enumerate,1]{leftmargin=1.5em,itemsep=1em}
\setlist[enumerate,2]{label=(\alph*),leftmargin=2em,itemsep=2.2em}

\newif\ifkey
\ifdefined\iskey \keytrue \fi
\colorlet{keyred}{red!75!black}
\newcommand{\ans}[1]{\ifkey\leavevmode{\color{keyred}#1}\fi}
% vertical work space on the student copy, less on the key
\newcommand{\work}[1]{\ifkey\vspace{0.6em}\else\vspace{#1}\fi}
\newcommand{\fig}[1]{\input{../../figures/#1}}
% Title line: \lessontitle{Quartiles, IQR, and Box Plots}
\newcommand{\lessontitle}[1]{{\Large\bfseries #1\ifkey\ \textcolor{keyred}{KEY}\fi}\hfill Name \underline{\hspace{2in}}\par\vspace{0.6em}}
\newcommand{\partner}{\par\vspace{0.4em}Work with a partner. Compare answers when you finish each problem.\par}


\begin{document}

\lessontitle{Simulations}

\renewcommand{\arraystretch}{1.3}
\begin{tabular}{|>{\raggedright\arraybackslash}p{1.6in}|>{\raggedright\arraybackslash}p{2.7in}|>{\raggedright\arraybackslash}p{2.1in}|}\hline
 & \textbf{How you find it} & \textbf{Coin example}\\ \hline
\textbf{Theoretical probability} & By reasoning about equally likely outcomes. & $P(\text{heads}) = \frac{1}{2}$\\[0.3em] \hline
\textbf{Experimental probability} & By doing trials:\par\rule{0pt}{2.2em}$\dfrac{\text{number of times it happened}}{\text{number of trials}}$ & 20 flips, 12 heads:\par $\frac{12}{20} = 0.6$\\[0.9em] \hline
\end{tabular}

\vspace{0.6em}
\textbf{The law of large numbers:} the more trials you do, the closer the experimental probability gets to the theoretical probability.

\vspace{1em}
\textbf{A simulation} uses random numbers to act out a chance event many times.

\vspace{0.3em}
\begin{tabular}{@{}ll}
1. & Assign digits to the outcomes so the probabilities match.\\
2. & Decide what one \textbf{trial} is.\\
3. & Run many trials with random digits.\\
4. & Count how often the event happened.\\
\end{tabular}

\vspace{0.8em}
\textbf{Example:} A basketball player makes 70\% of her free throws. About how often does she make all 3 of 3 shots?

\vspace{0.3em}
\begin{tabular}{@{}ll}
1. & Digits 0 to 6 = make (7 of the 10 digits is 70\%). \; Digits 7, 8, 9 = miss.\\
2. & One trial = 3 digits (3 shots).\\
3. & Use row 1 below. Make 10 trials.\\
\end{tabular}

\vspace{0.4em}
\begin{center}
\renewcommand{\arraystretch}{1.3}{\large
\begin{tabular}{|l|l|}\hline
Row 1 & \texttt{\bfseries 26647 37828 83005 96123 30669 70587}\\
Row 2 & \texttt{\bfseries 15021 10728 56722 07313 39118 78635}\\ \hline
\end{tabular}}
\end{center}

\begin{center}
\renewcommand{\arraystretch}{1.3}
\setlength{\tabcolsep}{4pt}
\begin{tabular}{|l|c|c|c|c|c|c|c|c|c|c|}\hline
Trial & 1 & 2 & 3 & 4 & 5 & 6 & 7 & 8 & 9 & 10 \\ \hline
Digits & 266 & 473 & 782 & 883 & 005 & 961 & 233 & 066 & 970 & 587 \\ \hline
Makes & \textbf{3} & 2 & 1 & 1 & \textbf{3} & 2 & \textbf{3} & \textbf{3} & 1 & 1 \\ \hline
\end{tabular}
\end{center}

4 of the 10 trials made all 3. \; Experimental probability: $\frac{4}{10} = 0.4$.

\vspace{0.2em}
The theoretical probability is $0.7 \cdot 0.7 \cdot 0.7 = 0.343$. With more trials, the simulation would get closer to it.

\partner

% ---------------- PROBLEMS ----------------
\newpage
\begin{enumerate}

\item A spinner has 4 equal parts: red, blue, green, yellow. It is spun 40 times.

\begin{center}
\renewcommand{\arraystretch}{1.3}
\begin{tabular}{|c|c|c|c|}\hline
Red & Blue & Green & Yellow \\ \hline
14 & 8 & 10 & 8 \\ \hline
\end{tabular}
\end{center}
\begin{enumerate}[itemsep=1.4em]
  \item Find the experimental probability of red. \ans{\quad $\frac{14}{40} = 0.35$}
  \item Find the theoretical probability of red. \ans{\quad $\frac{1}{4} = 0.25$}
  \item Why are they different? \ans{\quad Chance. 40 spins is not very many trials.}
  \item What would you expect if the spinner were spun 4,000 times? \ans{\quad The fraction of reds would be close to 0.25.}
\end{enumerate}
\work{1em}

\item How would you assign random digits to simulate each event?
\begin{enumerate}[itemsep=1.6em]
  \item A coin lands heads ($P = 0.5$). \ans{\quad 0 to 4 = heads, 5 to 9 = tails}
  \item A student is chosen from a class where 30\% walk to school. \ans{\quad 0, 1, 2 = walks; 3 to 9 = does not}
  \item A number cube is rolled. \ans{\quad 1 to 6 = the number rolled; skip 0, 7, 8, 9}
  \item An event with $P = 0.25$. \ans{\quad Use pairs: 00 to 24 = yes, 25 to 99 = no}
\end{enumerate}

\newpage
\item A family has 3 children. Each child is equally likely to be a girl or a boy.
Simulate: even digit = girl, odd digit = boy. One trial = 3 digits.
\begin{enumerate}[itemsep=1em]
  \item Use \textbf{row 2} on page 1 to run 10 trials. Count the girls in each trial.

  \begin{center}
  \renewcommand{\arraystretch}{1.5}
  \setlength{\tabcolsep}{4pt}
  \begin{tabular}{|l|p{0.38in}|p{0.38in}|p{0.38in}|p{0.38in}|p{0.38in}|p{0.38in}|p{0.38in}|p{0.38in}|p{0.38in}|p{0.38in}|}\hline
  Trial & 1 & 2 & 3 & 4 & 5 & 6 & 7 & 8 & 9 & 10 \\ \hline
  Digits & \ans{150} & \ans{211} & \ans{072} & \ans{856} & \ans{722} & \ans{073} & \ans{133} & \ans{911} & \ans{878} & \ans{635} \\ \hline
  Girls & \ans{1} & \ans{1} & \ans{2} & \ans{2} & \ans{2} & \ans{1} & \ans{0} & \ans{0} & \ans{2} & \ans{1} \\ \hline
  \end{tabular}
  \end{center}
  \item Estimate $P(\text{at least 2 girls})$ from your simulation. \ans{\quad $\frac{4}{10} = 0.4$}
  \item The theoretical probability is 0.5. How could you make your estimate closer? \ans{\quad Run many more trials.}
\end{enumerate}
\work{1.5em}

\item A student flips a coin 50 times and gets 34 heads. She wonders if the coin is fair.
A computer simulated 100 sets of 50 flips of a fair coin. The number of heads ranged from 17 to 33. No set had 34 or more heads.
\begin{enumerate}[itemsep=1em]
  \item What does the simulation tell you about getting 34 heads with a fair coin?

  \ans{It is very unusual. It did not happen once in 100 simulations.}
  \work{2em}
  \item Do you think the coin is fair? Explain.

  \ans{Probably not. 34 heads would almost never happen with a fair coin, so the result is statistically significant.}
\end{enumerate}

\end{enumerate}

\end{document}
